In unauthorized model distillation, distillers use outputs from proprietary models to train student models without the providers' permission.
We systematize existing defenses based on traffic screening, model watermarking, and passive post-hoc fingerprinting, and identify new limitations in these approaches.
A practical defense must combine high detection rates, strong evidence of training use, low operational cost, and robustness to data cleaning, but existing approaches only partially meet these goals.
We therefore study a setting in which providers coordinate defenses across the model-serving pipeline.
In this setting, we introduce \name{}, a behavioral watermarking method for detecting unauthorized distillation of language model agents.
\name{} inserts safe auxiliary tool calls into teacher trajectories to embed behaviors that students can learn during distillation.
The provider then uses traffic screening to select suspicious trajectories and targeted probes to test whether a suspect student has learned these behaviors.
We evaluate \name{} on three agent benchmarks with eight different teacher-student model pairs. 
It detects all 24 distilled students with no false positives in 48 control models.
Our experiments also show that \name{} preserves task utility and remains effective under data mixing, paraphrasing, truncation, and adaptive cleaning, demonstrating its practicality in real-world use scenarios.
